% =================================================
% 文件: 网络基础.tex
% 章节: 计算机网络基础
% 版本: v1.0
% =================================================

\documentclass[12pt,UTF8]{ctexbook}

\input{../common/page}

\xeCJKsetup{AutoFallBack=true}
\setCJKfamilyfont{hei}{SimHei}
\setCJKfamilyfont{kai}{KaiTi}

\usepackage{titletoc}
\titlecontents{chapter}[0pt]{\vspace{3mm}\bf\addvspace{2pt}\filright}{\contentspush{\thecontentslabel\hspace{0.8em}}}{}{\titlerule*[8pt]{.}\contentspage}

\usepackage[colorlinks,linkcolor=blue,citecolor=blue,urlcolor=blue]{hyperref}

\ctexset{
	part/name= {第,卷},
	part/number={\chinese{part}},
	chapter/name={第,章},
	chapter/number={\arabic{chapter}}
}

\usepackage{graphicx}
\graphicspath{{Images/}}

\usepackage[perpage]{footmisc}

\usepackage{enumitem}

\title{\heiti\zihao{0} 计算机网络基础入门指南}
\author{WangFei}
\date{\today}

\begin{document}

\maketitle
\tableofcontents

\frontmatter
\chapter{前言}

本指南面向初学者，详细介绍计算机网络的基础概念、TCP/IP协议、路由交换、网络安全、网络配置等知识。

\mainmatter

\chapter{网络基础概念}
\label{chap:network_basics}

\section{什么是计算机网络}
\label{sec:what_is_network}

计算机网络是将多台计算机通过通信设备连接起来，实现资源共享和信息交换的系统。

网络的分类：
\begin{table}[htbp]
    \centering
    \caption{网络分类}
    \begin{tabular}{|c|p{8cm}|}
        \hline
        \textbf{分类} & \textbf{说明} \\
        \hline
        LAN & 局域网，覆盖范围小 \\
        \hline
        WAN & 广域网，覆盖范围大 \\
        \hline
        MAN & 城域网，覆盖城市范围 \\
        \hline
        PAN & 个人局域网，覆盖个人设备 \\
        \hline
    \end{tabular}
\end{table}

\section{网络拓扑}
\label{sec:network_topology}

网络拓扑是网络中设备的连接方式：
\begin{table}[htbp]
    \centering
    \caption{网络拓扑}
    \begin{tabular}{|c|p{8cm}|}
        \hline
        \textbf{拓扑} & \textbf{说明} \\
        \hline
        总线型 & 所有设备连接到一条总线 \\
        \hline
        星型 & 所有设备连接到中心节点 \\
        \hline
        环型 & 设备连成环形 \\
        \hline
        网状 & 设备之间相互连接 \\
        \hline
    \end{tabular}
\end{table}

\section{OSI 七层模型}
\label{sec:osi_model}

OSI（Open Systems Interconnection）七层模型：

\begin{table}[htbp]
    \centering
    \caption{OSI 七层模型}
    \begin{tabular}{|c|c|p{6cm}|}
        \hline
        \textbf{层级} & \textbf{名称} & \textbf{功能} \\
        \hline
        7 & 应用层 & 提供应用服务 \\
        \hline
        6 & 表示层 & 数据格式转换 \\
        \hline
        5 & 会话层 & 建立和管理会话 \\
        \hline
        4 & 传输层 & 端到端传输 \\
        \hline
        3 & 网络层 & 路由和寻址 \\
        \hline
        2 & 数据链路层 & 帧的传输 \\
        \hline
        1 & 物理层 & 物理介质传输 \\
        \hline
    \end{tabular}
\end{table}

\chapter{TCP/IP 协议}
\label{chap:tcp_ip}

\section{TCP/IP 四层模型}
\label{sec:tcp_ip_model}

TCP/IP 四层模型：
\begin{table}[htbp]
    \centering
    \caption{TCP/IP 四层模型}
    \begin{tabular}{|c|c|p{6cm}|}
        \hline
        \textbf{层级} & \textbf{名称} & \textbf{协议} \\
        \hline
        4 & 应用层 & HTTP、FTP、DNS \\
        \hline
        3 & 传输层 & TCP、UDP \\
        \hline
        2 & 网络层 & IP、ICMP、ARP \\
        \hline
        1 & 网络接口层 & Ethernet、PPP \\
        \hline
    \end{tabular}
\end{table}

\section{IP 地址}
\label{sec:ip_address}

IP 地址分为 IPv4 和 IPv6：

IPv4 地址格式：32位二进制，通常表示为点分十进制

IP 地址分类：
\begin{table}[htbp]
    \centering
    \caption{IPv4 地址分类}
    \begin{tabular}{|c|c|c|}
        \hline
        \textbf{类别} & \textbf{范围} & \textbf{用途} \\
        \hline
        A类 & 1-126 & 大型网络 \\
        \hline
        B类 & 128-191 & 中型网络 \\
        \hline
        C类 & 192-223 & 小型网络 \\
        \hline
        D类 & 224-239 & 组播 \\
        \hline
        E类 & 240-255 & 保留 \\
        \hline
    \end{tabular}
\end{table}

子网掩码：用于划分网络和主机部分

CIDR 表示法：192.168.1.0/24

配置 IPv4：
\begin{verbatim}
nmcli connection modify eth0 ipv4.addresses 192.168.1.100/24
nmcli connection modify eth0 ipv4.gateway 192.168.1.1
nmcli connection modify eth0 ipv4.dns 192.168.1.1
nmcli connection modify eth0 ipv4.method manual
nmcli connection up eth0
\end{verbatim}

\section{TCP 协议}
\label{sec:tcp}

TCP（传输控制协议）特点：
\begin{itemize}
    \item 面向连接
    \item 可靠传输
    \item 有序传输
    \item 流量控制
    \item 拥塞控制
\end{itemize}

TCP 三次握手：
\begin{enumerate}
    \item SYN：客户端发送同步请求
    \item SYN+ACK：服务器响应并确认
    \item ACK：客户端确认
\end{enumerate}

TCP 四次挥手：
\begin{enumerate}
    \item FIN：客户端发送终止请求
    \item ACK：服务器确认
    \item FIN：服务器发送终止请求
    \item ACK：客户端确认
\end{enumerate}

\section{UDP 协议}
\label{sec:udp}

UDP（用户数据报协议）特点：
\begin{itemize}
    \item 无连接
    \item 不可靠传输
    \item 速度快
    \item 适合实时应用
\end{itemize}

UDP 适用场景：
\begin{itemize}
    \item 视频流
    \item 音频流
    \item 实时游戏
    \item DNS 查询
\end{itemize}

\chapter{路由与交换}
\label{chap:routing}

\section{路由基础}
\label{sec:routing_basics}

路由是数据包从源到目的地的过程。

路由类型：
\begin{table}[htbp]
    \centering
    \caption{路由类型}
    \begin{tabular}{|c|p{8cm}|}
        \hline
        \textbf{类型} & \textbf{说明} \\
        \hline
        静态路由 & 手动配置的路由 \\
        \hline
        动态路由 & 自动学习的路由 \\
        \hline
        默认路由 & 没有明确路由时使用 \\
        \hline
    \end{tabular}
\end{table}

动态路由协议：
\begin{table}[htbp]
    \centering
    \caption{动态路由协议}
    \begin{tabular}{|c|p{8cm}|}
        \hline
        \textbf{协议} & \textbf{说明} \\
        \hline
        RIP & 路由信息协议，距离向量 \\
        \hline
        OSPF & 开放最短路径优先，链路状态 \\
        \hline
        BGP & 边界网关协议，外部路由 \\
        \hline
    \end{tabular}
\end{table}

配置静态路由：
\begin{verbatim}
ip route add 10.0.0.0/8 via 192.168.1.254
ip route add default via 192.168.1.1
\end{verbatim}

查看路由表：
\begin{verbatim}
ip route show
route -n
\end{verbatim}

\section{交换基础}
\label{sec:switching}

交换机工作在数据链路层，根据 MAC 地址转发数据。

交换机类型：
\begin{table}[htbp]
    \centering
    \caption{交换机类型}
    \begin{tabular}{|c|p{8cm}|}
        \hline
        \textbf{类型} & \textbf{说明} \\
        \hline
        二层交换机 & 基于 MAC 地址转发 \\
        \hline
        三层交换机 & 支持路由功能 \\
        \hline
        管理型交换机 & 可配置管理 \\
        \hline
        非管理型交换机 & 即插即用 \\
        \hline
    \end{tabular}
\end{table}

VLAN（虚拟局域网）：
\begin{itemize}
    \item 将物理网络划分为多个逻辑网络
    \item 提高网络安全性
    \item 减少广播风暴
\end{itemize}

\chapter{网络设备}
\label{chap:network_devices}

\section{常见网络设备}
\label{sec:common_devices}

\begin{table}[htbp]
    \centering
    \caption{网络设备}
    \begin{tabular}{|c|p{8cm}|}
        \hline
        \textbf{设备} & \textbf{说明} \\
        \hline
        网卡 & 计算机网络接口 \\
        \hline
        集线器 & 物理层设备，共享带宽 \\
        \hline
        交换机 & 数据链路层设备 \\
        \hline
        路由器 & 网络层设备，连接不同网络 \\
        \hline
        防火墙 & 安全设备，过滤流量 \\
        \hline
    \end{tabular}
\end{table}

\section{网络接口配置}
\label{sec:interface_config}

查看网络接口：
\begin{verbatim}
ip addr show
ifconfig
\end{verbatim}

配置网络接口（CentOS Stream 9）：
\begin{verbatim}
nmcli connection add type ethernet con-name eth0 ifname eth0
nmcli connection modify eth0 ipv4.addresses 192.168.1.100/24
nmcli connection modify eth0 ipv4.gateway 192.168.1.1
nmcli connection modify eth0 ipv4.dns 8.8.8.8
nmcli connection up eth0
\end{verbatim}

\chapter{网络服务}
\label{chap:network_services}

\section{DNS}
\label{sec:dns}

DNS（域名系统）将域名解析为 IP 地址。

DNS 查询过程：
\begin{enumerate}
    \item 检查本地缓存
    \item 查询本地 DNS 服务器
    \item 查询根域名服务器
    \item 查询顶级域名服务器
    \item 查询权威域名服务器
\end{enumerate}

配置 DNS：
\begin{verbatim}
cat /etc/resolv.conf
nameserver 8.8.8.8
nameserver 8.8.4.4
\end{verbatim}

测试 DNS：
\begin{verbatim}
nslookup example.com
dig example.com
\end{verbatim}

\section{DHCP}
\label{sec:dhcp}

DHCP（动态主机配置协议）自动分配 IP 地址。

DHCP 工作流程：
\begin{enumerate}
    \item DHCP Discover：客户端发现服务器
    \item DHCP Offer：服务器提供地址
    \item DHCP Request：客户端请求地址
    \item DHCP Acknowledge：服务器确认
\end{enumerate}

配置 DHCP 客户端：
\begin{verbatim}
nmcli connection modify eth0 ipv4.method auto
nmcli connection up eth0
\end{verbatim}

\section{HTTP}
\label{sec:http}

HTTP（超文本传输协议）用于 Web 通信。

HTTP 请求方法：
\begin{table}[htbp]
    \centering
    \caption{HTTP 请求方法}
    \begin{tabular}{|c|p{8cm}|}
        \hline
        \textbf{方法} & \textbf{说明} \\
        \hline
        GET & 请求资源 \\
        \hline
        POST & 提交数据 \\
        \hline
        PUT & 更新资源 \\
        \hline
        DELETE & 删除资源 \\
        \hline
    \end{tabular}
\end{table}

HTTP 状态码：
\begin{table}[htbp]
    \centering
    \caption{HTTP 状态码}
    \begin{tabular}{|c|c|p{5cm}|}
        \hline
        \textbf{范围} & \textbf{类别} & \textbf{说明} \\
        \hline
        1xx & 信息响应 & 请求已接收 \\
        \hline
        2xx & 成功响应 & 请求成功 \\
        \hline
        3xx & 重定向 & 需要进一步操作 \\
        \hline
        4xx & 客户端错误 & 请求有问题 \\
        \hline
        5xx & 服务器错误 & 服务器有问题 \\
        \hline
    \end{tabular}
\end{table}

\chapter{网络安全}
\label{chap:network_security}

\section{网络安全威胁}
\label{sec:security_threats}

常见网络安全威胁：
\begin{itemize}
    \item \textbf{ARP 欺骗}：伪造 ARP 响应
    \item \textbf{DDoS 攻击}：分布式拒绝服务
    \item \textbf{中间人攻击}：窃取通信数据
    \item \textbf{端口扫描}：探测开放端口
\end{itemize}

\section{安全防护}
\label{sec:security_protection}

网络安全防护措施：
\begin{itemize}
    \item \textbf{防火墙}：过滤网络流量
    \item \textbf{加密}：保护数据传输
    \item \textbf{VPN}：安全远程访问
    \item \textbf{入侵检测}：监控异常行为
\end{itemize}

\chapter{网络故障排查}
\label{chap:troubleshooting}

\section{故障排查工具}
\label{sec:troubleshooting_tools}

常用网络工具：
\begin{table}[htbp]
    \centering
    \caption{网络工具}
    \begin{tabular}{|c|p{8cm}|}
        \hline
        \textbf{工具} & \textbf{说明} \\
        \hline
        ping & 测试连通性 \\
        \hline
        traceroute & 追踪路由路径 \\
        \hline
        netstat & 查看网络连接 \\
        \hline
        tcpdump & 抓包分析 \\
        \hline
        nmap & 端口扫描 \\
        \hline
    \end{tabular}
\end{table}

使用 ping：
\begin{verbatim}
ping -c 4 192.168.1.1
ping -c 4 example.com
\end{verbatim}

使用 traceroute：
\begin{verbatim}
traceroute example.com
mtr example.com
\end{verbatim}

使用 tcpdump：
\begin{verbatim}
tcpdump -i eth0 port 80
tcpdump -i eth0 host 192.168.1.100
\end{verbatim}

\section{故障排查流程}
\label{sec:troubleshooting_flow}

网络故障排查步骤：
\begin{enumerate}
    \item 确认故障现象
    \item 检查物理连接
    \item 检查网络配置
    \item 测试连通性
    \item 抓包分析
    \item 定位问题
\end{enumerate}

\chapter{IPv6}
\label{chap:ipv6}

\section{IPv6 基础}
\label{sec:ipv6_basics}

IPv6 地址格式：128位，冒号分隔的十六进制

IPv6 地址表示：
\begin{itemize}
    \item 压缩前导零：2001:0db8:0000:0000:0000:0000:0000:0001
    \item 压缩连续零：2001:db8::1
\end{itemize}

IPv6 地址类型：
\begin{table}[htbp]
    \centering
    \caption{IPv6 地址类型}
    \begin{tabular}{|c|p{8cm}|}
        \hline
        \textbf{类型} & \textbf{说明} \\
        \hline
        单播地址 & 单个接口地址 \\
        \hline
        组播地址 & 一组接口地址 \\
        \hline
        任播地址 & 最近接口地址 \\
        \hline
    \end{tabular}
\end{table}

配置 IPv6：
\begin{verbatim}
nmcli connection modify eth0 ipv6.addresses 2001:db8::1/64
nmcli connection modify eth0 ipv6.gateway 2001:db8::254
nmcli connection modify eth0 ipv6.method manual
nmcli connection up eth0
\end{verbatim}

\backmatter

\end{document}